\section{Integration Model and Architecture}\label{sec:arch}

\subsection{The integration model}

In xlog a neural network \emph{is} a predicate. A declaration binds a network to a predicate symbol; evaluating that predicate runs a forward pass, and the network's softmax over a label set can become a distribution over weighted facts. Typed device-buffer interfaces make those values available to the supported reasoning routes. Exact gradients of query values flow back through the probabilistic knowledge-compilation route; deterministic and epistemic execution reuse the integration surface without claiming that gradient semantics. Facts, circuit values, and gradient buffers remain device-resident at their data-plane handoffs, while orchestration and bounded control reads remain host responsibilities.
\subsection{GPU residency model}

xlog keeps runtime data-plane state in GPU memory and does not silently spill it or switch to a host executor. The provider owns one device handle, stream pool, asynchronous resource, runtime, and memory manager. Exceeding a configured byte budget raises \texttt{RESOURCE\_EXHAUSTED}; exhausting a fixed row, tuple, sample, or other count capacity raises \texttt{CAPACITY\_EXCEEDED}; and a bounded fixpoint that reaches its iteration limit raises \texttt{CONVERGENCE\_FAILURE}. Transfer counters delimit measured claims. In particular, the certified resident recursive and Monte Carlo sampled cores record zero tracked transfers, then return one bounded terminal receipt whose status is authoritative.

The payoff is bandwidth asymmetry. PCIe~4.0 $\times16$ delivers roughly 32\,GB/s, while GPU HBM provides 1--3\,TB/s. A single unnecessary D2H round-trip for a moderate relation (10M tuples at 16 bytes, ${\approx}160$\,MB) costs about 5\,ms. Keeping fact stores, deltas, circuit values, solver state, and gradient buffers on the GPU avoids copying those data planes through host memory; this claim does not erase the bounded control observations documented for each route.

\subsection{Crate architecture}

\begin{figure*}[tbp]
  \centering
  \begin{tikzpicture}[x=1mm,y=1mm,
      cell/.style={fnode,anchor=north west,font=\footnotesize,
                   inner xsep=2mm,inner ysep=1.6mm,minimum height=12mm},
      tag/.style={ftier,anchor=north west,text width=21mm,minimum height=12mm,
                  font=\footnotesize}]
    % ---- Tier 4 --- interfaces (host-facing) ----
    \node[tag] (t4tag) at (0,0)
      {\textbf{Tier 4}\\{\scriptsize\scshape interfaces}};
    \node[cell,text width=67.5mm] (t4c1) at (27,0)
      {Python API (PyO3)\\[0.6mm]\cratechip{pyxlog}};
    \node[cell,text width=67.5mm] (t4c2) at (100.5,0)
      {command-line interface\\[0.6mm]\cratechip{xlog-cli}};
    % ---- Tier 3 --- pipelines ----
    \node[tag] (t3tag) at (0,-19.5)
      {\textbf{Tier 3}\\{\scriptsize\scshape pipelines}};
    \node[cell,text width=43mm] (t3c1) at (27,-19.5)
      {deterministic\\[0.6mm]\cratechip{xlog-gpu}};
    \node[cell,text width=43mm] (t3c2) at (76,-19.5)
      {probabilistic $\cdot$ PIR $\cdot$ verified KC\\[0.6mm]\cratechip{xlog-prob}};
    \node[cell,text width=43mm] (t3c3) at (125,-19.5)
      {neural-symbolic\\[0.6mm]\cratechip{xlog-neural} \cratechip{xlog-induce}};
    % ---- Tier 2 --- subsystems ----
    \node[tag] (t2tag) at (0,-36.5)
      {\textbf{Tier 2}\\{\scriptsize\scshape subsystems}};
    \node[cell,text width=43mm] (t2c1) at (27,-36.5)
      {language frontend\\[0.6mm]\cratechip{xlog-logic}};
    \node[cell,text width=43mm] (t2c2) at (76,-36.5)
      {GPU runtime executor\\[0.6mm]\cratechip{xlog-runtime}};
    \node[cell,text width=43mm] (t2c3) at (125,-36.5)
      {GPU CDCL $\cdot$ local search\\[0.6mm]\cratechip{xlog-solve}};
    % ---- Tier 1 --- IR & device services ----
    \node[tag] (t1tag) at (0,-53.5)
      {\textbf{Tier 1}\\{\scriptsize\scshape ir \& device\\[-0.4mm]\scriptsize\scshape services}};
    \node[cell,text width=43mm] (t1c1) at (27,-53.5)
      {relational / epistemic IRs\\[0.6mm]\cratechip{xlog-ir}};
    \node[cell,text width=43mm] (t1c2) at (76,-53.5)
      {CUDA wrapper, Arrow/DLPack\\[0.6mm]\cratechip{xlog-cuda}};
    \node[cell,text width=43mm] (t1c3) at (125,-53.5)
      {runtime statistics\\[0.6mm]\cratechip{xlog-stats}};
    % ---- Tier 0 --- core ----
    \node[tag] (t0tag) at (0,-70.5)
      {\textbf{Tier 0}\\{\scriptsize\scshape core}};
    \node[cell,text width=141mm] (t0c1) at (27,-70.5)
      {scalar types $\cdot$ errors $\cdot$ symbol table $\cdot$ memory budgets\\[0.6mm]\cratechip{xlog-core}};
    % ---- GPU device plane behind tiers 1-3 ----
    \begin{scope}[on background layer]
      \node[fplane,fit=(t3tag)(t3c3)(t1tag)(t1c3),inner sep=2.2mm] (plane) {};
    \end{scope}
    \node[fplanetitle,anchor=east] at ([xshift=-4mm]plane.north east)
      {\planetitle{runtime data plane $\cdot$ gpu-resident}};
    % ---- dependency arrows (tier-level, down the tag column) ----
    \draw[farrow] (12.5,-12) -- (12.5,-19.5);
    \draw[farrow] (12.5,-31.5) -- (12.5,-36.5);
    \draw[farrow] (12.5,-48.5) -- (12.5,-53.5);
    \draw[farrow] (12.5,-65.5) -- (12.5,-70.5);
    \node[flabel,anchor=west] at (14,-14.2) {depends on};
    % ---- host annotations ----
    \node[flabel,anchor=south east] at (172,1.5)
      {host: orchestration $\cdot$ I/O $\cdot$ compilation};
    \node[flabel,anchor=north east] at (172,-84)
      {host-side foundation: shared types \& traits};
  \end{tikzpicture}
  \caption{xlog crate hierarchy. Each tier depends only on lower tiers; the
  shaded device plane marks GPU-resident runtime data planes, while the host
  provides orchestration, I/O, compilation, and bounded control reads.}
  \label{fig:arch}
\end{figure*}

The workspace is organized into five tiers with strict layering: no crate depends on a peer or a higher tier (\Cref{fig:arch}). A leaf \emph{core} crate defines shared scalar types, error and memory-budget types, canonical boolean parsing, and a bidirectional symbol table for dictionary-encoded strings. Above it, a tier of intermediate representations and device services wraps CUDA via cudarc~\cite{cudarc}, resolves embedded kernel modules, and exposes Arrow and DLPack interoperability. Each CUDA provider owns its complete runtime resource graph; there is no process-wide runtime singleton or alternate provider path. The three core subsystems (the typed language frontend of \Cref{sec:lang}, the GPU runtime executor, and the GPU CDCL/local-search solver) compose into integrated pipelines for deterministic, probabilistic, and neural-symbolic execution, which a PyO3~\cite{pyo3} extension module (\texttt{pyxlog}) and a command-line binary surface to users.

\subsection{From program to GPU plan}

A program is compiled to a GPU-executable plan in five stages.

\textbf{Parsing.} A PEG parser built with pest produces an AST covering the full surface, including probabilistic facts, annotated disjunctions, epistemic literals, and neural-predicate declarations.

\textbf{Dependency analysis} builds a graph with positive, negative, aggregate, probabilistic, and modal edges and computes its strongly connected components (SCCs) with Tarjan's algorithm. It establishes a deterministic stratum order for the ordinary monotone core and classifies supported nonmonotone components for their dedicated plans.

\textbf{Lowering} fans out from the normalized AST into three reasoning IRs: deterministic and reduced ordinary rules become relational-IR trees, probabilistic provenance becomes PIR, and modal programs retain their semantics in EIR. Neural predicates compose with the applicable route. SAT/MaxSAT and verification use a separate shared service instead of a fourth reasoning IR.

\textbf{Optimization} pushes predicates through joins, while lowering chooses a deterministic greedy join order. Shape-specific selectivity rewrites cover recognized triangle and four-cycle plans; general dynamic-programming join ordering is not active. Eligible multiway rules are promoted to the WCOJ subsystem, and bound recursive queries can be rewritten via magic sets (\Cref{sec:datalog}).

\textbf{Plan assembly} emits the route-specific execution plan, whose relational work is dispatched to CUDA kernels. Supported negated recursion uses GPU-backed well-founded semantics instead of the ordinary semi-naive schedule.

\subsection{Reasoning IRs and circuit format}

Three reasoning IRs keep the backends decoupled; probabilistic knowledge compilation also emits a downstream circuit format. Each form has one role.

\textbf{RIR (Relational IR)} is the algebraic tree (\texttt{Scan}, \texttt{Filter}, \texttt{Project}, \texttt{Join}, \texttt{GroupBy}, \texttt{Union}, \texttt{Distinct}, \texttt{Diff}, \texttt{Fixpoint}, and the worst-case-optimal \texttt{MultiWayJoin}/\texttt{ChainJoin}) carrying cardinality, skew, and delta-vs-full metadata. Deterministic execution and reduced ordinary subprograms consume RIR; the other backends preserve additional semantics in their own IRs.

\textbf{PIR (Provenance IR)} is the weighted Boolean formula (\texttt{Lit}, \texttt{NegLit}, \texttt{And}, \texttt{Or}, \texttt{Decision}) derived from provenance tracking over a probabilistic program, capturing its structure without execution order.

\textbf{XGCF (xlog GPU Circuit Format)} is the compiled, device-resident form: a levelized DAG whose level offsets let every node at one topological level evaluate in a single kernel launch, with forward (log-space weighted model counting) and backward (adjoint) passes.

\textbf{EIR (Epistemic IR)} preserves modal operators for world-view reasoning and lowers to a dedicated GPU execution plan instead of ordinary relational algebra (\Cref{sec:epistemic}).

\textbf{GpuCnf and CDCL} form a shared solver service, not a fourth reasoning IR. Callers provide a device-resident CNF formula for SAT/MaxSAT search or verification. The CDCL service consumes that formula directly; XGCF remains the circuit format produced only by probabilistic knowledge compilation.
